\scanpage{30}
\textbf{C. R. Acad. Sc. Paris, t. 292 (19 janvier 1981)\hfill Série I --- }\mspage{215}

\textsc{Théorème} 2 (cf. [3]). --- \textit{Soit} $m\in\mathbf{N}^{*}$ \textit{et soit} $\operatorname{Mon}(\mathbf{R}^m)$ \textit{le monoïde des applications continues et injectives de} $\mathbf{R}^m$ \textit{en} $\mathbf{R}^m$. \textit{Si} $\Omega$ \textit{est un faisceau sur} $\mathscr{V}_m$ [\textit{au sens de l'axiome} $(a)$] \textit{les conditions ci-dessous sont équivalentes} :

(i) $\Omega$ \textit{est localement constant} (cf. [4]);

(ii) $\Omega$ \textit{est séparé} [\textit{au sens de l'axiome} $(b)$];

(iii) \textit{Le monoïde} $\operatorname{Mon}(\mathbf{R}^m)$ \textit{opère sur} $\Omega(\mathbf{R}^m)$ \textit{par automorphismes}.

Soit $\mathrm{L.\ C.}\,(\mathscr{V}_m)$ la catégorie des faisceaux localement constants sur $\mathscr{V}_m$ et soit $\Pi(m)$ le groupe symétrisé du monoïde $\operatorname{Mon}(\mathbf{R}^m)$.

\textsc{Théorème} 3 (cf. [3]). --- \textit{Le foncteur} [cf. \textit{théorème} 2, (iii)] :
\[
\Omega\to\Omega(\mathbf{R}^m):\quad\mathrm{L.\ C.}\,(\mathscr{V}_m)\to\Pi(m)-\mathrm{Ens},
\]
\textit{est une équivalence de catégories} [\textit{i. e.} $\Pi_1(\tilde{\mathscr{V}}_m)\overset{\sim}{\to}\Pi(m)$].

La théorie du degré de Brouwer fourni un homomorphisme de groupes :
\[
\deg_m:\quad\Pi(m)\to\mu_2.
\]

\textsc{Théorème} 4 (cf. [3]). --- \textit{Les conditions ci-après sont équivalentes} :

$(a)$ \textit{l'homomorphisme} $\deg_m:\Pi(m)\overset{\sim}{\to}\mu_2$ \textit{est un isomorphisme};

$(b)$ \textit{tout homéomorphisme de} $\mathbf{R}^m$ \textit{de degré} $+1$ \textit{est « stable »} (\textit{i. e. composé fini d'homéomorphismes dont chacun est l'identité sur un ouvert non vide}).

Puisque Kirby a prouvé que pour $m\ne4$ tout homéomorphisme de $\mathbf{R}^m$ de degré $+1$ est stable (\textit{cf.} [2]) on a le résultat suivant :

\textsc{Corollaire}. --- \textit{Si} $m\ne4$ \textit{l'homomorphisme} $\deg_m::\Pi(m)\overset{\sim}{\to}\mu_2$ \textit{est un isomorphisme}.

\textsc{Définition}. --- On appelle \textit{triplet de Stokes} (en dimension $m$) tout triplet $\tau=(\Omega_m,\Omega_{m-1},\delta)$ d'un faisceau séparé $\Omega_m$ sur $\mathscr{V}_m$, d'un faisceau séparé $\Omega_{m-1}$ sur $\mathscr{V}_{m-1}$ et d'un morphisme fonctoriel :
\[
\delta:\quad\Omega_m^{\bullet}\to\Omega_{m-1}\circ\partial
\]
($\partial:\mathscr{V}_m\to\mathscr{V}_{m-1}$ est le foncteur bord) qui soit un « isomorphisme local » au sens de l'axiome $(c)$.

La notion de morphisme entre triplets de Stokes en dimension $m$ étant évidente, on note par $\operatorname{St}(m)$ la catégorie des triplets de Stokes en dimension $m$.

\textsc{Théorème} 5 (cf. [3]. --- \textit{Le foncteur} :
\[
(\Omega_m,\Omega_{m-1},\delta)\to\Omega_m(\mathbf{R}^m):\operatorname{St}(m)\overset{\approx}{\to}\Pi(m)-\mathrm{Ens},
\]
\textit{est une équivalence de catégories}.

(*) Remise le 5 janvier 1981.

[1] R. \textsc{Godement}, \textit{Théorie des faisceaux}, Herman, Paris, 1958.

[2] R. \textsc{Kirby}, \textit{Ann. of Math.}, 89, 1969.

[3] M. \textsc{Wanderley}, \textit{Théorie axiomatique de l'orientation des variétés topologiques à bord} (\textit{Thèse de} 3\textsuperscript{e} \textit{cycle}, Montpellier, 1980).

[4] SGA 4, 1\textsuperscript{er} vol., \textit{Séminaire de Géométrie algébrique de l'I.H.E.S.}, dirigé par M. \textsc{Artin}, A. \textsc{Grothendieck} et J. L. \textsc{Verdier} (\textit{Lecture Notes in Math.} 289, Springer).

\textit{U.E.R. de Mathématiques, Université des Sciences et Techniques du Languedoc\unsure{.|,|;|:}}\\
\textit{Ed. N. S. das Graças, App. 202, Rua das Graças, 277, 50.000 Recife, P. E., Bresil.}

C. R., 1981, 1\textsuperscript{er} \textit{Semestre}. (T. 292, N\textsuperscript{o} \textbf{3})\hfill\textbf{Série I --- 15}

% archive sheet number: 29
